\chapter{Algèbre~: équations et inéquations}\label{ch:g10:algebra}

L'algèbre est l'art de calculer avec des lettres. Ce chapitre revoit et
affine les deux gestes de base --- \omterm{def:g10:algebra:expand}{développer} et \omterm{def:g10:algebra:expand}{factoriser} --- et les
utilise pour résoudre systématiquement des \omterm{def:g10:algebra:equation}{équations} et des inéquations,
avec le \omterm{def:g10:algebra:signtable}{tableau de signes} comme outil central.

\section{Développer et factoriser}

\begin{definition}[Développement, factorisation]\label{def:g10:algebra:expand}
\emph{Développer}\index{développer} une expression, c'est transformer des
produits en sommes~; \emph{factoriser}\index{factoriser}, c'est
l'opération inverse, transformer une somme en un produit. La règle de
base est la distributivité~:
\[
k(a + b) = ka + kb,
\qquad
(a + b)(c + d) = ac + ad + bc + bd .
\]
\end{definition}

\begin{example}\label{ex:g10:algebra:expand}
Développer $(2x + 3)(x - 5)$ étape par étape~:
\begin{align*}
(2x + 3)(x - 5)
&= 2x \times x + 2x \times (-5) + 3 \times x + 3 \times (-5) \\
&= 2x^2 - 10x + 3x - 15 \\
&= 2x^2 - 7x - 15 .
\end{align*}
\end{example}

\begin{theorem}[Identités remarquables]\label{thm:g10:algebra:identities}
Pour tous \omterm{def:g10:numbers:sets}{nombres réels} $a$ et $b$~:
\[
(a+b)^2 = a^2 + 2ab + b^2,
\qquad
(a-b)^2 = a^2 - 2ab + b^2,
\qquad
(a+b)(a-b) = a^2 - b^2 .
\]
\end{theorem}

\begin{proof}
Chacune est un développement direct. Pour la première~:
\[
(a+b)^2 = (a+b)(a+b) = a^2 + ab + ba + b^2 = a^2 + 2ab + b^2 .
\]
La deuxième s'obtient en remplaçant $b$ par $-b$~:
$(a-b)^2 = a^2 + 2a(-b) + (-b)^2 = a^2 - 2ab + b^2$. Pour la troisième~:
$(a+b)(a-b) = a^2 - ab + ba - b^2 = a^2 - b^2$.
\end{proof}

\begin{omfigure}
\begin{tikzpicture}[scale=0.62]
  \draw[thick] (0,0) rectangle (5,5);
  \draw[thick] (3.4,0) -- (3.4,5);
  \draw[thick] (0,3.4) -- (5,3.4);
  \fill[omDef!12] (0,0) rectangle (3.4,3.4);
  \fill[omThm!12] (3.4,3.4) rectangle (5,5);
  \fill[omProp!12] (3.4,0) rectangle (5,3.4);
  \fill[omProp!12] (0,3.4) rectangle (3.4,5);
  \node[font=\small] at (1.7,1.7) {$a^2$};
  \node[font=\small] at (4.2,4.2) {$b^2$};
  \node[font=\small] at (4.2,1.7) {$ab$};
  \node[font=\small] at (1.7,4.2) {$ab$};
  \node[below, font=\small] at (1.7,0) {$a$};
  \node[below, font=\small] at (4.2,0) {$b$};
  \node[left, font=\small] at (0,1.7) {$a$};
  \node[left, font=\small] at (0,4.2) {$b$};
\end{tikzpicture}

{\small Pourquoi $(a+b)^2 = a^2 + 2ab + b^2$~: un carré de côté $a + b$ se
découpe en un carré de côté $a$, un carré de côté $b$, et deux rectangles
$a \times b$.}
\end{omfigure}

\begin{example}\label{ex:g10:algebra:identities}
Utilisées dans le sens direct (\omterm{def:g10:algebra:expand}{développer}) et inverse (\omterm{def:g10:algebra:expand}{factoriser})~:
\begin{align*}
(3x + 2)^2 &= 9x^2 + 12x + 4,
& x^2 - 25 &= (x+5)(x-5), \\
(x - 4)^2 &= x^2 - 8x + 16,
& 4x^2 + 4x + 1 &= (2x + 1)^2 .
\end{align*}
Le calcul mental en profite aussi~: $101 \times 99 = (100+1)(100-1) =
10000 - 1 = 9999$.
\end{example}

\begin{method}[Factoriser une expression]\label{met:g10:algebra:factor}
Essayer, dans cet ordre~:
\begin{enumerate}
  \item \emph{Facteur commun~:} repérer un facteur présent dans chaque
        terme et le sortir, par exemple $6x^2 + 4x = 2x(3x + 2)$. Le
        facteur commun peut être une parenthèse entière~:
        $(x+1)(2x-3) + (x+1)(x+7) = (x+1)(3x+4)$.
  \item \emph{Identité remarquable~:} reconnaître $a^2 - b^2$, ou
        $a^2 \pm 2ab + b^2$, par exemple $9x^2 - 16 = (3x)^2 - 4^2 =
        (3x+4)(3x-4)$.
  \item Combiner les deux, et vérifier le résultat en le redéveloppant.
\end{enumerate}
\end{method}

\begin{example}\label{ex:g10:algebra:factorsteps}
Factoriser $E = (2x + 1)^2 - (x - 3)^2$. C'est une différence de deux
carrés $a^2 - b^2$ avec $a = 2x+1$ et $b = x-3$~:
\begin{align*}
E &= \bigl[(2x+1) + (x-3)\bigr] \times \bigl[(2x+1) - (x-3)\bigr] \\
  &= (3x - 2)(2x + 1 - x + 3) \\
  &= (3x - 2)(x + 4).
\end{align*}
Attention à la seconde parenthèse~: soustraire $x - 3$, c'est soustraire
$x$ \emph{et} ajouter $3$.
\end{example}

\section{Équations}

\begin{definition}[Équation, solution]\label{def:g10:algebra:equation}
Une \emph{équation}\index{équation} est une égalité contenant un nombre
inconnu $x$~; une \emph{solution} est une valeur de $x$ qui rend l'égalité
vraie. Résoudre l'équation, c'est trouver \emph{toutes} ses solutions.
\end{definition}

Deux \omterm{def:g10:algebra:equation}{équations} sont \emph{équivalentes} lorsqu'elles ont exactement les
mêmes solutions. Ajouter le même nombre aux deux membres, ou multiplier
les deux membres par le même nombre \emph{non nul}, produit une \omterm{def:g10:algebra:equation}{équation}
équivalente.

\begin{example}[Équation du premier degré]\label{ex:g10:algebra:firstdegree}
Résoudre $5x - 7 = 2x + 8$, une opération à la fois~:
\begin{align*}
5x - 7 &= 2x + 8 \\
5x - 2x &= 8 + 7 && \text{(ajouter $7 - 2x$ aux deux membres)} \\
3x &= 15 \\
x &= 5 && \text{(diviser les deux membres par $3$).}
\end{align*}
L'unique solution est $5$. Vérification~: $5 \times 5 - 7 = 18$ et
$2 \times 5 + 8 = 18$.
\end{example}

\begin{theorem}[Règle du produit nul]\label{thm:g10:algebra:zeroproduct}
Un produit de \omterm{def:g10:numbers:sets}{nombres réels} est nul si et seulement si l'un au moins de
ses facteurs est nul~:
\[
A \times B = 0
\quad\Longleftrightarrow\quad
A = 0 \ \text{ ou } \ B = 0 .
\]
\end{theorem}

\begin{proof}
Si $A = 0$ ou $B = 0$, le produit est clairement $0$. Réciproquement,
supposons $AB = 0$ et $A \neq 0$. Diviser les deux membres par $A$ donne
$B = 0$.
\end{proof}

\begin{method}[Résoudre par factorisation]\label{met:g10:algebra:zeroproduct}
Pour résoudre une \omterm{def:g10:algebra:equation}{équation} dont le second membre est $0$ et dont le
premier membre peut se \omterm{def:g10:algebra:expand}{factoriser}~:
\begin{enumerate}
  \item tout ramener à gauche pour que l'\omterm{def:g10:algebra:equation}{équation} s'écrive $E = 0$~;
  \item \omterm{def:g10:algebra:expand}{factoriser} $E$ en un produit de facteurs du premier degré~;
  \item appliquer la règle du produit nul et résoudre chaque petite
        \omterm{def:g10:algebra:equation}{équation}~;
  \item rassembler toutes les solutions trouvées.
\end{enumerate}
\end{method}

\begin{example}\label{ex:g10:algebra:zeroproduct}
Résoudre $(x - 2)(3x + 6) = 0$~: soit $x - 2 = 0$, d'où $x = 2$, soit
$3x + 6 = 0$, d'où $x = -2$. Solutions~: $-2$ et $2$.

Résoudre $x^2 = 9x$. \emph{Ne pas} diviser par $x$ (il pourrait être
nul~!)~; tout ramener d'un côté et \omterm{def:g10:algebra:expand}{factoriser}~:
\[
x^2 - 9x = 0
\quad\Longleftrightarrow\quad
x(x - 9) = 0
\quad\Longleftrightarrow\quad
x = 0 \ \text{ ou } \ x = 9 .
\]
\end{example}

\begin{example}[Équation quotient]\label{ex:g10:algebra:quotient}
Résoudre $\dfrac{x - 1}{x + 2} = 0$. Un quotient est nul exactement lorsque
son numérateur est nul \emph{et} son dénominateur ne l'est pas. Ici
$x - 1 = 0$ donne $x = 1$, et $1 + 2 = 3 \neq 0$~: l'unique solution est
$1$. La valeur $x = -2$ est \emph{interdite} (division par zéro) et doit
toujours être exclue en premier.
\end{example}

\section{Inéquations et tableaux de signes}

\begin{proposition}[Règles pour les inégalités]\label{prop:g10:algebra:ineqrules}
Soit $a \leq b$.
\begin{enumerate}
  \item Pour tout $c$~: $a + c \leq b + c$ (ajouter préserve l'ordre).
  \item Si $c > 0$~: $ac \leq bc$ (multiplier par un nombre positif
        préserve l'ordre).
  \item Si $c < 0$~: $ac \geq bc$ (multiplier par un nombre négatif
        \emph{inverse} l'ordre).
\end{enumerate}
\end{proposition}

\begin{proof}
1.~$(b + c) - (a + c) = b - a \geq 0$.
2.~$bc - ac = (b-a)c$ est le produit de deux nombres positifs ou nuls,
donc $bc \geq ac$.
3.~À présent $(b-a)c \leq 0$ comme produit d'un nombre positif ou nul et
d'un nombre négatif, donc $bc \leq ac$.
\end{proof}

\begin{example}\label{ex:g10:algebra:firstineq}
Résoudre $-2x + 3 > 9$~:
\begin{align*}
-2x + 3 &> 9 \\
-2x &> 6 && \text{(soustraire $3$)} \\
x &< -3 && \text{(diviser par $-2 < 0$~: inverser l'inégalité).}
\end{align*}
L'ensemble des solutions est $\intoo{-\infty}{-3}$.
\end{example}

\begin{definition}[Tableau de signes]\label{def:g10:algebra:signtable}
Un \emph{tableau de signes}\index{tableau de signes} enregistre, une
ligne par facteur, le signe ($+$ ou $-$) de chaque facteur d'une
expression lorsque $x$ parcourt $\R$, avec un $0$ à chaque valeur où le
facteur s'annule~; une dernière ligne donne le signe du produit, en
appliquant la règle des signes colonne par colonne.
\end{definition}

\begin{example}\label{ex:g10:algebra:signtable}
Signe de $P(x) = (x - 2)(3 - x)$. Le facteur $x - 2$ s'annule en $2$ et
est positif après~; le facteur $3 - x$ s'annule en $3$ et est positif
avant~:
\begin{center}
\renewcommand{\arraystretch}{1.3}
\begin{tabular}{c|ccccc}
$x$ & $-\infty$ \quad & $2$ & & $3$ & \quad $+\infty$ \\
\hline
$x - 2$ & $-$ & $0$ & $+$ & & $+$ \\
$3 - x$ & $+$ & & $+$ & $0$ & $-$ \\
\hline
$P(x)$ & $-$ & $0$ & $+$ & $0$ & $-$ \\
\end{tabular}
\end{center}
Donc $P(x) > 0$ exactement sur $\intoo{2}{3}$, et $P(x) \leq 0$ sur
$\intoc{-\infty}{2} \cup \intco{3}{+\infty}$.
\end{example}

\begin{omfigure}
\begin{tikzpicture}
\begin{axis}[omaxis,
  width=8.2cm, height=5.2cm,
  xmin=-0.6, xmax=5.4, ymin=-2.6, ymax=1.4,
  xtick={2,3}, ytick=\empty]
  \addplot[omDef, thick, domain=0.35:4.65] {(x-2)*(3-x)};
  \fill (axis cs:2,0) circle (1.5pt);
  \fill (axis cs:3,0) circle (1.5pt);
  \node[omDef, font=\small] at (axis cs:2.5,0.65) {$P > 0$};
  \node[omThm, font=\small] at (axis cs:0.9,-1.3) {$P < 0$};
  \node[omThm, font=\small] at (axis cs:4.1,-1.3) {$P < 0$};
\end{axis}
\end{tikzpicture}

{\small Le graphique de $P(x) = (x-2)(3-x)$ confirme le \omterm{def:g10:algebra:signtable}{tableau de
signes}~: positif entre les racines $2$ et $3$, négatif à l'extérieur.}
\end{omfigure}

\begin{method}[Résoudre une inéquation avec un tableau de signes]\label{met:g10:algebra:signtable}
Pour résoudre une inéquation du type $E(x) \geq 0$ ou $E(x) < 0$~:
\begin{enumerate}
  \item tout ramener au premier membre, pour que le second membre soit
        $0$ (ne jamais multiplier les deux membres par une expression dont
        le signe est inconnu)~;
  \item \omterm{def:g10:algebra:expand}{factoriser} le premier membre, ou le mettre sur un dénominateur
        commun s'il s'agit d'un quotient~;
  \item construire le \omterm{def:g10:algebra:signtable}{tableau de signes} avec une ligne par facteur (pour
        un quotient, marquer les valeurs interdites d'une double barre
        dans la dernière ligne)~;
  \item lire les \omterm{def:g10:numbers:interval}{intervalles} où le signe demandé est réalisé, en
        vérifiant si chaque borne est incluse.
\end{enumerate}
\end{method}

\begin{example}\label{ex:g10:algebra:quotientineq}
Résoudre $\dfrac{x + 1}{x - 3} \leq 0$. Le numérateur s'annule en $-1$, le
dénominateur en $3$ (valeur interdite). Le quotient est négatif lorsque
les deux ont des signes opposés, ce qui arrive pour $x$ entre $-1$ et $3$.
En incluant la borne où le \emph{numérateur} s'annule mais en excluant la
valeur interdite~: l'ensemble des solutions est $\intco{-1}{3}$.
\end{example}

\section{Exercices}

\begin{exercise}[$\star$]\label{exo:g10:algebra:1}
Développer et simplifier~:
\[
3(2x - 5) - 2(x - 7), \qquad
(x + 4)(2x - 1), \qquad
(3x - 2)^2, \qquad
(5 - x)(5 + x).
\]
\end{exercise}

\begin{exercise}[$\star$]\label{exo:g10:algebra:2}
Factoriser~:
\[
x^2 - 49, \qquad
x^2 + 10x + 25, \qquad
7x^2 - 14x, \qquad
16x^2 - 8x + 1 .
\]
\end{exercise}

\begin{exercise}[$\star$]\label{exo:g10:algebra:3}
Résoudre~:
\[
4x + 9 = x - 3, \qquad
\frac{2x - 1}{3} = 5, \qquad
2(x - 3) = 2x + 1 .
\]
(L'une d'elles n'a pas de solution --- expliquer pourquoi.)
\end{exercise}

\begin{exercise}[$\star$]\label{exo:g10:algebra:4}
Résoudre en utilisant la règle du produit nul~:
\[
(x + 5)(2x - 8) = 0, \qquad
x^2 - 16 = 0, \qquad
x^2 + 6x + 9 = 0 .
\]
\end{exercise}

\begin{exercise}[$\star$]\label{exo:g10:algebra:5}
Résoudre les inéquations et écrire les ensembles de solutions comme
\omterm{def:g10:numbers:interval}{intervalles}~:
\[
3x - 5 \leq 7, \qquad
-4x + 1 > 9, \qquad
2(x - 1) \geq 5x + 4 .
\]
\end{exercise}

\begin{exercise}[$\star\star$]\label{exo:g10:algebra:6}
Factoriser $E = (x - 1)(x + 3) - (x - 1)(2x - 5)$, puis résoudre
$E = 0$.
\end{exercise}

\begin{exercise}[$\star\star$]\label{exo:g10:algebra:7}
Résoudre $x^2 = 5x$, puis $(2x + 1)^2 = (x - 4)^2$. (Pour la seconde,
tout ramener à gauche et \omterm{def:g10:algebra:expand}{factoriser} la différence de carrés.)
\end{exercise}

\begin{exercise}[$\star\star$]\label{exo:g10:algebra:8}
À l'aide d'un \omterm{def:g10:algebra:signtable}{tableau de signes}, résoudre
\[
(x + 2)(4 - x) > 0,
\qquad
\frac{2x - 6}{x + 1} \geq 0 .
\]
\end{exercise}

\begin{exercise}[$\star\star$]\label{exo:g10:algebra:9}
Un service de streaming coûte $9$ par mois, plus des frais d'inscription
uniques de $15$~; un concurrent coûte $12$ par mois sans frais. Au bout de
combien de mois le premier service devient-il le moins cher au total~?
Mettre en \omterm{def:g10:algebra:equation}{équation} et résoudre une inéquation.
\end{exercise}

\begin{exercise}[$\star\star$]\label{exo:g10:algebra:10}
Résoudre $\dfrac{x+3}{x-1} = 2$. (Multiplier les deux membres par $x - 1$
après avoir exclu la valeur interdite, ou tout ramener à gauche sur un
dénominateur commun.)
\end{exercise}

\begin{exercise}[$\star\star\star$]\label{exo:g10:algebra:11}
Soit $n$ un entier. Montrer que $(n + 1)^2 - n^2$ est toujours impair, et
que la différence des carrés de deux nombres impairs consécutifs est
toujours un multiple de $8$.
\end{exercise}

\section{Problème~: de Babylone au nombre d'or}

\begin{problem}[{Devoir du week-end --- problèmes somme-produit,
la mise sous forme canonique telle qu'al-Khwarizmi la dessinait, et
l'équation $x^2 = x + 1$}]
\label{pb:g10:algebra:1}
Il y a quatre mille ans, les scribes babyloniens résolvaient ceci sur
l'argile~: \emph{trouver deux nombres connaissant leur somme et leur
produit}. Leur astuce --- chercher les nombres symétriquement autour
de la demi-somme --- vient encore aujourd'hui à bout de toute \omterm{def:g10:algebra:equation}{équation}
du second degré, sous son nom moderne~: la \emph{mise sous forme
canonique}. Ce problème apprend l'astuce, regarde al-Khwarizmi
compléter littéralement un carré, puis met la méthode au travail sur
la plus célèbre de toutes les proportions.

\medskip
\noindent\textbf{Partie I --- Le problème babylonien.}
\begin{enumerate}
  \item Échauffement (\cref{thm:g10:algebra:identities})~: \omterm{def:g10:algebra:expand}{factoriser}
        $x^2 - 10x + 25$, puis $9x^2 - 49$, puis --- en combinant une
        identité et une différence de carrés --- $x^2 + 6x + 9 - 4$.
  \item Résoudre $(x + 1)(x + 5) = 0$ et $(3x - 7)(3x + 7) = 0$
        (\cref{thm:g10:algebra:zeroproduct}).
  \item Le problème du scribe~: deux nombres ont pour somme $10$ et
        pour produit $21$. À la manière de Babylone, les écrire
        $5 - t$ et $5 + t$ (symétriques autour de la demi-somme) et
        déterminer $t$, puis les deux nombres.
  \item Même méthode~: somme $14$, produit $33$.
  \item La recette générale~: pour une somme $s$ et un produit $p$,
        montrer que les deux nombres valent
        \[
        \frac s2 - t
        \quad\text{et}\quad
        \frac s2 + t,
        \qquad\text{où}\quad
        t^2 = \left(\frac s2\right)^2 - p ,
        \]
        et énoncer la condition sur $s$ et $p$ pour que le problème
        ait une solution.
\end{enumerate}

\medskip
\noindent\textbf{Partie II --- La mise sous forme canonique.}
\begin{enumerate}[resume]
  \item Développer $(x - a)(x - b)$ et conclure~: les solutions de
        $x^2 - sx + p = 0$ sont exactement les couples de nombres de
        somme $s$ et de produit $p$. (Les scribes résolvaient des
        \omterm{def:g10:algebra:equation}{équations} du second degré sans le savoir.)
  \item Vérifier que $x^2 - 10x + 21 = (x - 5)^2 - 4$, puis résoudre
        $x^2 - 10x + 21 = 0$ en factorisant le membre de droite comme
        une différence de carrés. Comparer avec la question 3.
  \item Résoudre $x^2 + 6x - 7 = 0$ par mise sous forme canonique.
  \item Mettre $x^2 + 4x + 5$ sous forme canonique, et expliquer
        pourquoi l'\omterm{def:g10:algebra:equation}{équation} $x^2 + 4x + 5 = 0$ n'a \emph{aucune}
        solution réelle. Énoncer le critère général~: une fois écrite
        $(x + h)^2 + k$, quand une expression du second degré
        s'annule-t-elle~?
  \item Le nom est à prendre au pied de la lettre. À Bagdad, au
        neuvième siècle, al-Khwarizmi dessinait $x^2 + 6x = 7$ comme
        un carré de côté $x$ auquel on colle sur deux côtés deux
        rectangles $3 \times x$, puis il \emph{complétait le carré}
        par le coin $3 \times 3$ manquant. Faire la figure et
        expliquer comment elle transforme l'\omterm{def:g10:algebra:equation}{équation} en
        $(x + 3)^2 = 16$ --- puis résoudre.
\end{enumerate}

\medskip
\noindent\textbf{Partie III --- Signes, tableaux et deux vieilles
promesses.}
\begin{enumerate}[resume]
  \item Construire le \omterm{def:g10:algebra:signtable}{tableau de signes} de $(x - 3)(x + 2)$
        (\cref{def:g10:algebra:signtable}) et résoudre
        $(x - 3)(x + 2) < 0$.
  \item Résoudre $x^2 \leq 10x - 21$ (tout ramener dans un même
        membre et réutiliser la question 7).
  \item Résoudre $\dfrac{x - 1}{x + 2} \geq 0$ à l'aide d'un \omterm{def:g10:algebra:signtable}{tableau
        de signes}, sans oublier la valeur interdite.
  \item Une vieille promesse tenue~: pour l'agriculteur du problème de
        la clôture de la reine Didon, dans le volume précédent, avec
        $20$ m de clôture, l'aire d'un enclos $x \times (10 - x)$
        vérifie
        \[
        A(x) = x(10 - x) = 25 - (x - 5)^2 .
        \]
        Vérifier l'identité et y lire une \emph{démonstration} de ce
        que l'on ne pouvait alors que conjecturer~: l'aire maximale,
        et l'unique $x$ qui la réalise.
  \item Une autre~: démontrer que $x + \dfrac1x \geq 2$ pour tout
        $x > 0$, avec égalité seulement en $x = 1$. (Multiplier par
        $x$, regrouper, reconnaître un carré --- l'inégalité
        arithmético-géométrique du devoir du week-end sur les
        relations d'Euclide, du volume précédent, en habit neuf.)
\end{enumerate}

\medskip
\noindent\textbf{Partie IV --- La section dorée.}
Découper un segment de longueur $1$ en une grande partie $x$ et une
petite partie $1 - x$ de sorte que
\emph{le tout soit à la grande partie ce que la grande partie est à la
petite}~: $\frac{1}{x} = \frac{x}{1 - x}$.
\begin{enumerate}[resume]
  \item Montrer que la condition s'écrit $x^2 + x - 1 = 0$.
  \item Résoudre par mise sous forme canonique, et conserver la
        racine qui est une longueur~:
        $x = \frac{\sqrt5 - 1}{2} \approx 0.618$.
  \item Le \emph{nombre d'or} est
        $\varphi = \frac1x = \frac{2}{\sqrt5 - 1}$. Rendre le
        dénominateur rationnel (multiplier haut et bas par
        $\sqrt5 + 1$~: une identité à l'œuvre) et montrer que
        $\varphi = \frac{\sqrt5 + 1}{2} \approx 1.618$.
  \item Vérifier \emph{exactement} les deux identités magiques de
        $\varphi$~:
        \[
        \varphi^2 = \varphi + 1,
        \qquad
        \frac{1}{\varphi} = \varphi - 1 .
        \]
  \item Pour finir~: expliquer pourquoi $\varphi$ est irrationnel
        (le devoir du week-end sur l'irrationalité, du volume
        précédent, fournit le fait clé sur $\sqrt5$), puis calculer à
        trois décimales les rapports $\frac53$, $\frac85$,
        $\frac{13}{8}$, $\frac{21}{13}$ de nombres de
        Hemachandra--Fibonacci consécutifs (devoir du week-end sur le
        décompte des rythmes, du volume précédent). Que semblent-ils
        poursuivre~? (La démonstration qu'ils l'atteignent relève des
        chapitres sur les suites.)
\end{enumerate}
\end{problem}
